% !TEX root = ../eigenfaces_algorithm.tex
%=====================================================
\section{收敛分析怎么呈现（对应题目的可衡量产出）}

画一张散点图：横轴 $|\lambda_{i+1}/\lambda_i|$，纵轴第 $i$ 对的迭代次数（对数轴）。
理论预测
\[
k_i\approx\frac{\ln(1/\varepsilon)}{\ln(\lambda_i/\lambda_{i+1})}
\]
是一条确定的曲线，把它画上去与实测点对照：
\begin{itemize}
  \item \textbf{吻合段}：谱间隙主导，幂迭代行为与理论一致；
  \item \textbf{偏离段}：从某个 $i$ 开始实测迭代数高于理论线——这正是收缩误差累积开始压过谱间隙效应的位置，
  两张图（迭代数 vs 间隙、相对误差 vs $i$）互为印证。
\end{itemize}
这张图是本题目最有“研究味”的产出：它把\textbf{课程里的谱半径收敛理论}
（Jacobi 迭代 $\rho(B)<1$ 的类比）迁移到了特征值算法上。
